\section{CLI 表示层：从用户意图到 Pipeline 契约}

前一章已经把 \texttt{pipeline.cpp} 解释为应用层聚合根（aggregate root）：它连接顺序读入、并行计算、
重排序缓冲区与单 writer。本章再向外走一层，讨论 \texttt{main.cpp} 这个命令行接口
（command-line interface, CLI）表示层（presentation layer）。它的任务不是“顺手启动程序”，
而是把用户在终端中写下的字符串序列
\[
    argv=(a_0,a_1,\ldots,a_{n-1})
\]
约束成一份可验证、可记录、可交给 pipeline 的执行契约。

与项目早期相比，当前 CLI 已经明显更像一个小型协议层：它负责格式推断、默认输出路径生成、覆盖保护、
\texttt{dry-run} 配置展示、布局转置策略与自动调优开关、文本/JSON 双报告，以及参数错误与执行错误的退出码分离。
因此，本章不再按“函数枚举”的方式展开，而是集中回答六个问题：表示层到底要解决什么需求，我们为何这样设计，用户到底可以说哪些命令，
命令如何流入 pipeline，哪些工程细节值得强调，以及当前边界在哪里。

\subsection{设计目标与表示层定位}

若把整个系统的依赖方向写成
\[
    \text{CLI}
    \longrightarrow
    \text{Pipeline}
    \longrightarrow
    \{\text{I/O},\text{Kernel}\}
    \longrightarrow
    \text{OS / CUDA Runtime},
\]
那么 CLI 层只拥有第一条箭头。它知道用户给了哪些路径、哪些数值参数、是否要覆盖已有文件、是否要输出 JSON；
但它不应该知道 \texttt{*.mat} 的内存映射细节，不应该读取矩阵元素，也不应该运行布局转置基准。表示层真正拥有的对象只有
“用户意图”，因此它最合理的定位就是一个意图编译器（intent compiler）：
\[
    \text{user intent}
    \longrightarrow
    \texttt{CliOptions}
    \longrightarrow
    \texttt{PipelineConfig}.
\]

从当前代码可反推，CLI 至少要满足四类互不重叠的需求。第一，它必须提供足够完整的表达能力，让用户能同时描述数据边界、
数值参数、布局策略、运行控制与报告方式。第二，它必须把高频路径压缩成短命令，例如只给输入文件时自动推导输出名，
否则实验脚本会充满重复样板。第三，它必须在进入昂贵的 I/O 与 GPU 计算之前尽可能早地拒绝错误，把字符串错误固定在表示层，
而不是把异常拖到 pipeline 或 kernel。第四，它必须是稳定契约，因为一旦文档示例、批处理脚本与论文实验开始依赖这些参数，
CLI 就成为事实上的用户空间接口。

这四类需求共同解释了当前实现为何采用“表驱动解析 + 归一化 + 验证 + 显式映射”的结构。单纯把
\texttt{argc/argv} 直接塞给下层当然也能“跑起来”，但那样会让表示层失去边界：路径推断散落在执行代码里，
错误信息靠异常偶遇，默认值无法审计，实验参数也难以复现。对于一个 CUDA 数值系统，这种混乱不是入口小问题，而是整个工程语义会漂移的起点。

\subsection{命令空间与选项契约}

当前 \texttt{ArgParser} 不是一串散落的 \texttt{if/else}，而是由 \texttt{OptionDefinition}
定义表驱动。长选项查找、短选项查找、是否需要值、帮助文本生成，全部共享同一张表。因此 CLI 的命令面并不是文档外加代码的
双份描述，而是由数据结构直接定义的单一事实来源（single source of truth）。这点很关键，因为一旦接口扩展，最怕的不是选项变多，
而是帮助文本、解析逻辑、默认语义三处开始不同步。

\begingroup
    \small
    \setlength{\tabcolsep}{4pt}
    \renewcommand{\arraystretch}{1.12}
    \begin{xltabular}{\linewidth}{@{}>{\raggedright\arraybackslash}p{0.12\linewidth}>{\raggedright\arraybackslash}p{0.29\linewidth}>{\raggedright\arraybackslash}p{0.12\linewidth}>{\raggedright\arraybackslash}X@{}}
        \caption{当前 CLI 命令面与选项契约}
        \label{tab:cli-command-surface}\\
        \toprule
        类别 & 选项或形式 & 参数 & 语义 \\
        \midrule
        \endfirsthead

        \multicolumn{4}{@{}l}{表~\thetable（续）当前 CLI 命令面与选项契约}\\
        \toprule
        类别 & 选项或形式 & 参数 & 语义 \\
        \midrule
        \endhead

        \midrule
        \multicolumn{4}{r@{}}{续下页}\\
        \endfoot

        \bottomrule
        \endlastfoot

        命令形式
        & \texttt{<input> [output]}\newline
          \texttt{--}
        & 路径或无
        & \texttt{<input>} 必填；\texttt{[output]} 省略时默认写入输入目录，文件名为 \texttt{<input-stem>.svd.<ext>}。
          显式 \texttt{--input}/\texttt{--output} 优先；\texttt{--} 之后 token 全按位置参数处理；多余位置参数报错。 \\

        数据边界
        & \texttt{--input}/\texttt{-i}\newline
          \texttt{--output}/\texttt{-o}
        & \texttt{PATH}
        & 输入路径必须存在且为文件；输出路径不能是目录，且规范化后不得与输入路径相同。
          输出已存在时必须显式给出 \texttt{--force}/\texttt{-y}。 \\

        格式契约
        & \texttt{--input-format}\newline
          \texttt{--output-format}\newline
          \texttt{--format}/\texttt{-f}
        & \texttt{FMT}
        & \texttt{FMT=auto|mat|txt}。\texttt{--format} 同时设置输入与输出格式。
          输入格式为 \texttt{auto} 且扩展名不可推断时直接报错；输出格式为 \texttt{auto} 时按“显式输入格式 $\rightarrow$ 输入扩展名 $\rightarrow$ \texttt{mat}”回退，并在缺失扩展名时自动补齐标准后缀。 \\

        数值配置
        & \texttt{--epsilon}/\texttt{-e}\newline
          \texttt{--max-sweeps}/\texttt{-s}\newline
          \texttt{--threads-per-block}/\texttt{-t}
        & \texttt{NUM}/\texttt{N}
        & 三项都要求正数（其中 \texttt{max-sweeps} 与 \texttt{threads-per-block} 为正整数）。
          默认值分别为 \texttt{1e-9}、\texttt{128}、\texttt{256}。 \\

        布局策略
        & \texttt{--layout-transpose-}\newline
          \texttt{mode|min-cols|}\newline
          \texttt{min-elems}
        & \texttt{MODE}/\texttt{N}
        & 选项族表达（family notation）：按“前缀 + 后缀”展开为完整选项。
          \texttt{MODE=auto|on|off}。
          两个阈值都要求正数，主要用于 \texttt{auto} 模式决策；默认值为 \texttt{auto}、\texttt{16}、\texttt{4096}。 \\

        自动调优
        & \texttt{--layout-transpose-}\newline
          \texttt{auto-tune|}\newline
          \texttt{bench-reps|}\newline
          \texttt{bench-sweeps}
        & 无或 \texttt{N}
        & 选项族表达（family notation）：按“前缀 + 后缀”展开为完整选项。
          开启后在 pipeline 前运行微基准自动调阈值；
          \texttt{bench-reps}/\texttt{bench-sweeps} 要求正整数，默认 \texttt{2}/\texttt{8}。
          仅当 \texttt{mode=auto} 时该调优实际生效，并用推荐值覆盖本次运行阈值。 \\

        运行控制
        & \texttt{--queue-capacity}/\texttt{-c}\newline
          \texttt{--force}/\texttt{-y}\newline
          \texttt{--dry-run}\newline
          \texttt{--print-config}
        & \texttt{N} 或无
        & \texttt{queue-capacity} 要求正整数，默认 \texttt{4}；\texttt{--force} 允许覆盖已有输出。
          \texttt{--dry-run} 只做参数校验并打印有效配置后退出（不执行 pipeline）；\texttt{--print-config} 会先打印配置再继续执行。 \\

        报告与动作
        & \texttt{--json-report}\newline
          \texttt{--quiet}/\texttt{-q}\newline
          \texttt{--help}/\texttt{-h}\newline
          \texttt{--version}/\texttt{-v}
        & 无
        & \texttt{--quiet} 仅抑制文本报告；\texttt{--json-report} 额外输出 JSON，二者可组合得到纯机器可读输出。
          \texttt{--help}/\texttt{--version} 属于动作命令：直接输出并以退出码 \texttt{0} 返回，不进入运行配置校验与 pipeline 执行。 \\
    \end{xltabular}
\endgroup


表~\ref{tab:cli-command-surface} 给出了当前 CLI 的完整命令面。这里有三点需要特别指出。第一，位置参数
\texttt{<input> [output]} 只是显式 \texttt{--input} 与 \texttt{--output} 的简写形式，因此
\[
    \texttt{jacobi-svd-cuda a.mat b.mat}
    \equiv
    \texttt{jacobi-svd-cuda --input a.mat --output b.mat}.
\]
第二，\texttt{--help} 与 \texttt{--version} 是动作，不是运行配置；它们请求的是“输出帮助/版本后退出”，而不是构造一个特殊 pipeline。
第三，布局转置相关参数被完整暴露出来，说明 CLI 不只是功能开关面板，也承担实验参数记录的职责。对 CUDA 程序而言，
\texttt{--layout-transpose-mode}、阈值和自动调优开关都属于可复现状态，而不是可有可无的 UI 细节。

\subsection{执行链路：从 \texttt{argv} 到 \texttt{run\_pipeline}}

当前控制流可以压缩为一条清晰的转换链：
\[
    \begin{aligned}
        argv
         & \xrightarrow{\text{token scan}}
        \texttt{ParseResult}
        \xrightarrow{\text{normalize}}
        \texttt{CliOptions}                \\
         & \xrightarrow{\text{validate}}
        \texttt{CliOptions}
        \xrightarrow{\Phi}
        \texttt{PipelineConfig}
        \xrightarrow{\text{run}}
        \texttt{PipelineReport}.
    \end{aligned}
\]

若展开为顺序流程，当前 CLI 大致经历六步。
\begin{enumerate}
    \item \textbf{词法扫描。} \texttt{ArgParser::parse} 对 \texttt{argv} 做一次线性扫描，把 token 分成长选项、短选项、位置参数与
          \texttt{--} 终止符。长选项支持 \texttt{--name=value} 与 \texttt{--name value}；短选项支持
          \texttt{-qv} 这样的无值聚合，也支持 \texttt{-e1e-10} 或 \texttt{-e 1e-10} 这样的带值形式。
    \item \textbf{语义归约。} 无论 token 来自长选项还是短选项，最终都经由 \texttt{apply\_option} 归约到同一份
          \texttt{ParseResult}，从而保证 \texttt{-i} 与 \texttt{--input} 只对应一种内部语义。
    \item \textbf{动作分流。} 解析结果先落入
          \[
              \texttt{ParseResult}=(\texttt{ParseAction},\texttt{CliOptions}),
          \]
          其中 $\texttt{ParseAction}\in\{\texttt{run},\texttt{help},\texttt{version}\}$。只有 \texttt{run}
          才进入归一化与验证；帮助和版本动作直接输出并成功返回。
    \item \textbf{归一化。} \texttt{normalize\_run\_options} 负责把不完整命令补成稳定配置：若输入格式为
          \texttt{auto}，则尝试从扩展名推断；若输出路径为空，则生成
          \[
              \texttt{<input-parent>/<input-stem>.svd.<ext>};
          \]
          若输出格式仍未显式给出，则用输入格式或扩展名推导 fallback，并在必要时补上规范扩展名。
    \item \textbf{运行前验证。} \texttt{validate\_run\_options} 检查输入是否存在、输入和输出是否为目录、输入输出是否规范化后相同，
          以及输出文件已存在时是否显式给出 \texttt{--force}。这一步的目标是把所有“本不该进入 pipeline 的错误”挡在入口。
    \item \textbf{执行与报告。} \texttt{main} 先截获 \texttt{dry-run} 和 \texttt{print-config}，随后把
          \texttt{CliOptions} 映射为 \texttt{PipelineConfig}，调用 \texttt{run\_pipeline}，最后按
          \texttt{quiet} 与 \texttt{json-report} 的组合输出文本摘要或 JSON 报告。
\end{enumerate}

从复杂度看，这条流程非常便宜。若 token 数为 $N$、选项定义数为 $M$，则解析阶段的查找代价近似为
\[
    O(NM).
\]
由于当前 $M$ 很小，实际表现接近线性。入口层的成本因此几乎可以忽略，真正昂贵的仍然是文件 I/O、CUDA kernel 和结果写回。

\subsection{关键工程细节}

\paragraph{第一，表驱动解析保证接口集中。}
\texttt{OptionDefinition} 同时承载内部标识、长名、短名、是否需要参数、值占位符与帮助描述。这样做的价值不是少写几行代码，
而是把“CLI 到底向用户承诺了什么”收敛到一份可审计数据。对后续维护而言，这远好于把命令含义散落在解析逻辑、帮助文本和错误消息里。

\paragraph{第二，字符串不会泄漏到应用层。}
当前实现把类型转换压缩进若干小函数，把自由字符串收缩为封闭集合
（closed set）、拒绝部分解析、非正数和非有限值。这样一来，下游看到的不是“也许合法的字符串”，而是已经满足前置条件
（precondition）的强类型参数。

\paragraph{第三，默认值服务高频路径，但不篡改显式意图。}
\texttt{CliOptions} 的默认值把最常见运行压缩成一条短命令：
\[
    \begin{aligned}
        \epsilon                     & = 10^{-9}, &
        \texttt{max\_sweeps}         & = 128,       \\
        \texttt{threads\_per\_block} & = 256,     &
        \texttt{queue\_capacity}     & = 4,
    \end{aligned}
\]
\[
    \begin{aligned}
        \texttt{layout\_transpose\_mode}          & = \texttt{auto}, \\
        \texttt{layout\_transpose\_min\_columns}  & = 16,            \\
        \texttt{layout\_transpose\_min\_elements} & = 4096.
    \end{aligned}
\]
自动调优默认关闭，\texttt{bench-reps} 默认为 $2$，\texttt{bench-sweeps} 默认为 $8$。若用户只给输入文件，
CLI 自动生成 \texttt{.svd} 输出名；但若用户已经显式给出输出路径，表示层不会粗暴重命名，只会在格式仍未明确时做最小必要补全。
这类细节体现的是工程边界感：系统应该替用户省去机械重复，而不是替用户改写明确意图。

\paragraph{第四，表示层控制位不会泄漏到 pipeline。}
\texttt{make\_pipeline\_config} 只把输入、输出、格式、队列窗口与 kernel 配置映射到
\texttt{PipelineConfig}。像 \texttt{dry\_run}、\texttt{print\_config}、\texttt{json\_report}、
\texttt{quiet}、\texttt{force\_overwrite} 这类纯表示层控制位，都在 \texttt{main.cpp} 中被消费。
这意味着下层不需要知道“用户是不是正在 dry-run”，只需要处理一份已经准备好的执行配置。这样的边界使应用层保持纯净，
也让测试更容易定位。

\paragraph{第五，错误分类和报告通道是显式契约。}
当前退出码被分成三类：
\[
    0:\text{成功，或 help/version/dry-run 等非执行动作成功},
\]
\[
    2:\text{CLI 参数错误},
    \qquad
    1:\text{执行失败}.
\]
这使脚本可以区分“命令写错了”和“命令合法但运行失败了”。成功路径上，文本摘要与 JSON 报告也是两条解耦通道：
\[
    \texttt{--quiet --json-report}
    \Rightarrow
    \text{只输出机器可读报告}.
\]
对实验自动化而言，这种清晰语义远比“反正 stdout 会打印点什么”更可靠。

\subsection{当前实现的边界与可演化方向}

当前 CLI 已经具备完整可用性，但边界也非常明确。第一，输入输出同一性判断仍基于
\texttt{absolute + lexically\_normal} 的词法规范化（path normalization）。这能防住
\texttt{a.mat} 与 \texttt{./a.mat} 之类的常见误用，却不能完整识别硬链接、符号链接或平台路径别名。
若未来进入更严苛环境，可以升级为平台文件身份比较，例如 POSIX 的 \texttt{st\_dev/st\_ino} 或 Windows file index。

第二，JSON 报告当前仍由手写 \texttt{std::cout} 直接拼接。由于现有字段几乎全是数值与固定枚举，
风险尚可接受；但如果未来把路径、设备名称、错误消息或驱动字符串纳入 JSON，就必须补上最小字符串转义、
稳定 schema 与版本字段。机器接口一旦被脚本依赖，就不应建立在“这个字段大概不会含引号”的侥幸上。

第三，当前命令来源只有 CLI 本身。这对可复现性是好事，因为一条命令就足以完整记录一次实验状态；
但若未来加入配置文件，优先级必须清晰写成
\[
    \text{default}\prec \text{config file}\prec \text{CLI override},
\]
并保留 \texttt{dry-run} 展示最终合并结果。否则，配置来源一多，实验者最容易失去对“这次到底按什么参数执行”的把握。

第四，\texttt{definitions\_} 这张选项表还没有被进一步复用。实际上，它天然适合生成 shell 补全
（shell completion）、接口文档，甚至参数兼容性检查。若后续命令面继续增长，这张表就是最自然的扩展支点。
同理，当前解析阶段对长名和短名采用线性查找是完全可接受的；只有当命令数继续增长，或者未来真的引入子命令
（subcommand）体系时，才值得把它演化为索引结构或分层命令模型。

第五，当前报告仍然是摘要型报告，而不是完整遥测（telemetry）。它告诉我们 testcase 数、总 sweep 数、
布局转置运行模式与总耗时，却没有每阶段耗时、队列阻塞时间、重排序缓冲区峰值或 GPU kernel 明细。若未来要做严肃 profiling，
更合理的方向是增设独立的遥测结构，而不是不断给稳定用户报告追加调试字段。

\subsection{本节小结}

这一章的核心是说明 \texttt{main.cpp} 如何建立一个边界清楚的表示层。用户写下的
\texttt{argv} 先被解析为 \texttt{ParseResult}，再经由归一化与验证收缩为 \texttt{CliOptions}，
随后映射为 \texttt{PipelineConfig}，最终得到 \texttt{PipelineReport} 与退出码。整条路径的价值在于：
字符串错误被尽早拒绝，默认值被集中管理，覆盖风险被显式化，实验参数可以被完整记录，而 pipeline 不必背负任何本属于表示层的偶然复杂度。

因此，CLI 在这个项目中不是外围装饰，而是系统契约的第一道强类型边界。它决定用户意图是否能被正确表达，
决定错误是在入口被发现还是拖到 GPU 上爆炸，也决定实验脚本、文档示例与自动化流程能否长期稳定运行。对于一个 CUDA 数值程序，
这类安静的入口代码和高吞吐 kernel 一样，都属于工程质量本身。
